\section*{Oppgave 2 - Produksjonsstyring}
\setcounter{section}{2}
\addcontentsline{toc}{section}{Oppgave 2 — Produksjonsstyring}

\setcounter{subsection}{0}
\subsection{Relevant teori}
Kundeordrens dekoblingspunkt (KODP) er det punktet i materialflyten som skiller produksjon basert på kundeordre fra produksjon basert på prognose eller lagernivå \parencite{strandhagen2024}. Plasseringen av KODP gir opphav til fire hovedstyringsformer: engineering på ordre (ETO), produksjon på ordre (MTO), montasje på ordre (ATO) og produksjon til lager (MTS). Når KODP ligger langt oppstrøms, slik som ved MTO, skjer en større del av produksjonen etter kundeordre. Det gir høy fleksibilitet, men lengre leveringstid. Når KODP ligger langt nedstrøms, slik som ved MTS, produseres mer på prognose, noe som gir kort leveringstid men øker risikoen for feil lagerbeholdning \parencite{slack2019}.

I et trykksystem (push) styres produksjonen av prognoser, mens et sugsystem (pull) lar faktisk forbruk nedstrøms utløse påfyll \parencite{strandhagen2024}. Lean-filosofien vektlegger å redusere sløsing, blant annet overproduksjon, unødvendig lager og transport \parencite{strandhagen2024}.


\subsection{Nåsituasjon hos Kvile}

Kvile bruker to ulike styringsformer. Vegvesenet styres som MTO, der produksjonen starter etter mottatt ordre. Butikksegmentet styres som MTS, der ferdig monterte og pakkede produkter ligger på ferdigvarelager. Forskjellene er oppsummert i tabell~\ref{tab:styringsformer} og illustrert i figur~\ref{fig:kodp}.

Begge segmentene benytter i praksis et trykksystem (push): for Vegvesenet trykkes produksjonen gjennom fabrikken av kundeordren, mens butikksegmentet trykkes av salgsprognoser. Produksjonen foregår i serier/batcher, der farge dikterer seriestørrelsen da betydelig omstillingstid på lakkanleggene gjør det nødvendig å samle opp produksjon av samme farge.

Sett i et Lean-perspektiv har Kviles nåværende oppsett potensielt flere former for sløsing: \textit{overproduksjon} av fargekombinasjoner med lavt salg (MTS med 304 SKU-er), \textit{unødvendig lager} bundet i ferdigvarer, \textit{venting} grunnet to døgns tørketid etter lakkering, og \textit{transport} med 10 lastebiler i tre-dagers turer gjennom hele Norge.

% --- Tabell: Styringsformer ---
\begin{table}[htbp]
    \centering
    \caption[Styringsformer]{Sammenligning av styringsformer hos Kvile AS.}
    \label{tab:styringsformer}
    \small
    \begin{tabular}{@{}lll@{}}
        \toprule
        & \textbf{Vegvesenet} & \textbf{Plantasjen/Møbelringen} \\
        \midrule
        Styringsform        & MTO                & MTS \\
        KODP                & Råvarelager         & Ferdigvarelager \\
        Styringsprinsipp    & Push (kundeordre)   & Push (prognose) \\
        Antall varianter    & 4                   & 304 \\
        Leveringstid        & $\geq$ 4 måneder   & Fra lager \\
        Planleggingshorisont & Lang, forutsigbar  & Kort, usikker \\
        \bottomrule
    \end{tabular}
\end{table}


\subsection{Vurdering og endringsforslag}

MTO-styringen av Vegvesenet er godt tilpasset både kundens og Kviles behov. Lang ordrehorisont, lavt antall varianter og forutsigbar etterspørsel gjør at Kvile kan planlegge kapasiteten effektivt uten å binde kapital i ferdigvarer. Denne styringsformen bør videreføres.

MTS-styringen for butikkmarkedet er derimot utfordrende. Med 304 varianter på ferdigvarelager blir kapitalbindingen stor, og risikoen for å sitte igjen med fargekombinasjoner som ikke selger er reell. Sesongsvingningene forsterker problemet — 50\,\% av salget skjer i mai--juli, noe som krever enten stor lageroppbygging i forkant eller kraftig kapasitetsøkning om våren. I tillegg vokser variantantallet med 96 nye SKU-er for hver ny modell.

\textbf{Forslag: Overgang til montasje på ordre (ATO).} Kvile kan flytte KODP fra ferdigvarelageret til komponentlageret ved å lagre ferdige, lakkerte halvfabrikata slik som aluminiumsrammer og trekomponenter separat. Når en ordre fra Plantasjen eller Møbelringen mottas, monteres riktig kombinasjon av ramme og tredeler, pakkes og klargjøres for henting. Dette er praktisk gjennomførbart fordi sluttmontasjen ifølge casebeskrivelsen kun krever enkle verktøy som skrutrekker og skiftenøkkel, og kjedene har tilstrekkelig ledetid (Plantasjen bestiller annenhver uke, Møbelringen månedlig).

Effekten på lagerkompleksiteten er stor. Tabell~\ref{tab:ato-sammenligning} viser hvordan antall lagervarer reduseres ved en overgang fra MTS til ATO.

% --- Tabell: MTS vs ATO lagervarer ---
\begin{table}[htbp]
    \centering
    \caption[MTS vs.\ ATO lagervarer]{Sammenligning av antall lagervarer (SKU-er) ved MTS og ATO for butikkmarkedet.}
    \label{tab:ato-sammenligning}
    \small
    \begin{tabular}{@{}lcc@{}}
        \toprule
        & \textbf{MTS (nå)} & \textbf{ATO (forslag)} \\
        \midrule
        Alu-rammer (4 modeller $\times$ 6 farger $\times$ delsett)
            & --- & $\sim$48 \\
        Trekomponenter (8 farger $\times$ delsett)
            & --- & $\sim$24 \\
        Ferdigvarer (bord + benk, alle kombinasjoner)
            & 304 & 0 \\
        \midrule
        \textbf{Totalt på lager} & \textbf{304} & $\mathbf{\sim}$\textbf{72} \\
        \textbf{Reduksjon} & & \textbf{$\sim$76\,\%} \\
        \bottomrule
    \end{tabular}
\end{table}

I tillegg bør Kvile vurdere å redusere farger med lav omløpshastighet og innføre syklisk styring av lakkanleggene.

% FIGUR: Kundeordrens dekoblingspunkt (KODP) hos Kvile AS


\begin{figure}[htbp]
    \centering
    \begin{tikzpicture}[
        node distance=0.5cm and 0.55cm,
        box/.style={draw, rounded corners=2pt, minimum height=0.75cm, minimum width=1.4cm, font=\scriptsize, align=center},
        progbox/.style={box, fill=gray!20},
        ordrebox/.style={box, fill=orange!20},
        arrw/.style={-{Stealth[length=2mm]}, semithick},
        lbl/.style={font=\tiny, text=black!55},
        >={Stealth[length=2mm]},
    ]

    % ===== VEGVESENET (MTO) =====
    \node[font=\scriptsize\bfseries, anchor=east] at (-0.3, 0) {Vegvesenet};
    \node[font=\tiny, anchor=east, text=black!50] at (-0.3, -0.3) {(MTO)};

    \node[progbox]  (v-inn)   at (1.0, 0) {Innkjøp};
    \node[ordrebox, right=of v-inn]  (v-prod) {Produksjon};
    \node[ordrebox, right=of v-prod] (v-mont) {Montasje};
    \node[ordrebox, right=of v-mont] (v-lev)  {Levering};
    \node[box, right=of v-lev, fill=blue!10]  (v-kunde) {Kunde};

    \draw[arrw] (v-inn) -- (v-prod);
    \draw[arrw] (v-prod) -- (v-mont);
    \draw[arrw] (v-mont) -- (v-lev);
    \draw[arrw] (v-lev) -- (v-kunde);

    % KODP
    \coordinate (v-kodp) at ($(v-inn.east)!0.5!(v-prod.west)$);
    \draw[red!70!black, thick, dashed] (v-kodp) ++(0, 0.55) -- ++(0, -1.1);
    \node[font=\tiny\bfseries, red!70!black, above] at ($(v-kodp)+(0, 0.55)$) {KODP};

    \node[lbl] at ($(v-inn)+(0, 0.6)$) {Prognose};
    \node[lbl] at ($(v-prod)+(0, 0.6)$) {Kundeordre};


    % ===== BUTIKK NÅ (MTS) =====
    \node[font=\scriptsize\bfseries, anchor=east] at (-0.3, -2.2) {Butikk \textit{(nå)}};
    \node[font=\tiny, anchor=east, text=black!50] at (-0.3, -2.5) {(MTS)};

    \node[progbox] (b-inn)   at (1.0, -2.2) {Innkjøp};
    \node[progbox, right=of b-inn]  (b-prod)  {Produksjon};
    \node[progbox, right=of b-prod] (b-mont)  {Montasje};
    \node[progbox, right=of b-mont] (b-lager) {FV-lager};
    \node[box, right=of b-lager, fill=blue!10] (b-kunde) {Kunde};

    \draw[arrw] (b-inn)  -- (b-prod);
    \draw[arrw] (b-prod) -- (b-mont);
    \draw[arrw] (b-mont) -- (b-lager);
    \draw[arrw] (b-lager) -- (b-kunde);

    % KODP
    \coordinate (b-kodp) at ($(b-lager.east)!0.5!(b-kunde.west)$);
    \draw[red!70!black, thick, dashed] (b-kodp) ++(0, 0.55) -- ++(0, -1.1);
    \node[font=\tiny\bfseries, red!70!black, above] at ($(b-kodp)+(0, 0.55)$) {KODP};

    \node[lbl] at ($(b-inn)+(0, 0.6)$) {Prognose};
    \node[lbl] at ($(b-mont)+(0, 0.6)$) {Prognose};

    % 304 SKU-er
    \node[font=\tiny, text=red!60!black, below=2pt of b-lager] {304 SKU-er};


    % ===== BUTIKK FORSLAG (ATO) =====
    \node[font=\scriptsize\bfseries, anchor=east] at (-0.3, -4.4) {Butikk \textit{(forslag)}};
    \node[font=\tiny, anchor=east, text=black!50] at (-0.3, -4.7) {(ATO)};

    \node[progbox] (a-inn)   at (1.0, -4.4) {Innkjøp};
    \node[progbox, right=of a-inn]  (a-prod)  {Prod.\ +\\[-2pt]lakk};
    \node[progbox, right=of a-prod] (a-lager) {Komp.-\\[-2pt]lager};
    \node[ordrebox, right=of a-lager] (a-mont) {Mont.\ +\\[-2pt]pakk};
    \node[box, right=of a-mont, fill=blue!10] (a-kunde) {Kunde};

    \draw[arrw] (a-inn)   -- (a-prod);
    \draw[arrw] (a-prod)  -- (a-lager);
    \draw[arrw] (a-lager) -- (a-mont);
    \draw[arrw] (a-mont)  -- (a-kunde);

    % KODP
    \coordinate (a-kodp) at ($(a-lager.east)!0.5!(a-mont.west)$);
    \draw[red!70!black, thick, dashed] (a-kodp) ++(0, 0.55) -- ++(0, -1.1);
    \node[font=\tiny\bfseries, red!70!black, above] at ($(a-kodp)+(0, 0.55)$) {KODP};

    \node[lbl] at ($(a-inn)+(0, 0.6)$) {Prognose};
    \node[lbl] at ($(a-mont)+(0, 0.6)$) {Kundeordre};

    % ~72 SKU-er
    \node[font=\tiny, text=green!50!black, below=2pt of a-lager] {$\sim$72 SKU-er};


    % ===== FORKLARING =====
    \node[draw, rounded corners, fill=white, inner sep=4pt, font=\tiny, anchor=north west] 
        at (-0.3, -5.8) {%
            \textcolor{gray!50}{\rule{0.4cm}{4pt}} Prognosestyrt \quad
            \textcolor{orange!45}{\rule{0.4cm}{4pt}} Ordrestyrt \quad
            \textcolor{red!70!black}{- - -} KODP
        };

    \end{tikzpicture}
    \caption[KODP for Kviles kundesegmenter]{Kundeordrens dekoblingspunkt (KODP) for Kviles to kundesegmenter. Øverst: Vegvesenet (MTO). Midten: butikkmarkedet i dag (MTS) med 304 SKU-er på ferdigvarelager. Nederst: foreslått ATO-løsning med $\sim$72 komponent-SKU-er.}
    \label{fig:kodp}
\end{figure}